\begin{Solbook}{1}
 \par \noindent  a) $-\dfrac {3}{10}$.\quad b) $\dfrac {41}{6}$.\quad c) $0$.\quad d) $\dfrac {2^{78}}{5\cdot 3^{26}}$.  \hfill \qedEX 
\end{Solbook}
\begin{Solbook}{2}
 \par \noindent  a) $x=-\dfrac {54}{25}$.\quad b) $x=\dfrac {1}{12}$.\quad c) $x=\dfrac {16}{29}$ (vì $x\neq 0$).\quad d) $x=\dfrac {3}{5}$ hoặc $x=\dfrac {19}{15}$.  \hfill \qedEX 
\end{Solbook}
\begin{Solbook}{3}
 \par \noindent  Người mua gạo tẻ trả ít hơn $4\%$.\\ Gọi giá và khối lượng gạo nếp lần lượt là $a$ và $m$. Tiền gạo nếp là $am$, tiền gạo tẻ là $0{,}8a\cdot 1{,}2m=0{,}96am$.  \hfill \qedEX 
\end{Solbook}
\begin{Solbook}{4}
 \par \noindent  $432\,\mathrm {m}^{3}$.\\ Sau lần một còn $\dfrac {2}{3}$ bể. Sau lần hai còn $\dfrac {1}{3}\cdot \dfrac {2}{3}=\dfrac {2}{9}$ bể. Gọi thể tích bể là $x$ ($\mathrm {m}^{3}$), $x>0$, ta có $\dfrac {2}{9}x=96$.  \hfill \qedEX 
\end{Solbook}
\begin{Solbook}{5}
 \par \noindent  $m\parallel n$.\\ Qua $H$ kẻ đường thẳng $d\parallel m$. Góc giữa $d$ và $AH$ bằng $40^{\circ }$. Vì $\widehat {AHB}=90^{\circ }$ nên góc giữa $d$ và $HB$ bằng $50^{\circ }$, bằng $\widehat {B_{1}}$. Suy ra $d\parallel n$. Mà $d\parallel m$ nên $m\parallel n$.  \hfill \qedEX 
\end{Solbook}
\begin{Solbook}{6}
 \par \noindent  a) $\widehat {ABC}=\widehat {ACB}=50^{\circ }$.\\ b) $AB=AC$, $BM=CM$, $AM$ chung nên $\triangle ABM=\triangle ACM$ (c.c.c). Hai góc tại $M$ bằng nhau và kề bù nên mỗi góc bằng $90^{\circ }$. Vậy $AM\perp BC$.\\ c) $M$ là trung điểm của $AD$ và $BC\perp AD$ nên $BC$ là đường trung trực của $AD$, do đó $DB=AB$. Tam giác $ABD$ cân tại $B$ và $\widehat {ADB}=40^{\circ }$.\\ \textit {Hướng dẫn ý c.} Chứng minh $\triangle ABM=\triangle DBM$ (c.g.c) để có $\widehat {BDM}=\widehat {BAM}=40^{\circ }$.  \hfill \qedEX 
\end{Solbook}
\begin{Solbook}{7}
 \par \noindent  $S=\dfrac {12}{25}$.\\ \textit {Hướng dẫn.} $\dfrac {1}{n(n+2)}=\dfrac {1}{2}\left (\dfrac {1}{n}-\dfrac {1}{n+2}\right )$. Tổng các số hạng rút gọn còn $\dfrac {1}{2}\left (1-\dfrac {1}{25}\right )$.  \hfill \qedEX 
\end{Solbook}
